\documentclass[11pt]{article}
\usepackage[margin=1in]{geometry}
\usepackage[T1]{fontenc}
\usepackage{lmodern}
\usepackage{amsmath}
\usepackage{enumitem}
\usepackage[colorlinks=true,linkcolor=blue,citecolor=blue,urlcolor=blue]{hyperref}
\usepackage[nameinlink,noabbrev]{cleveref}

\title{Deeply Nested Lists and Sectioning Stress Paper}
\author{Stage Two Corpus}
\date{}

\begin{document}
\maketitle

\begin{abstract}
This fixture exercises list and sectioning structures that frequently lose their
hierarchy during conversion. It contains a three-level nesting of an enumerate
inside an itemize inside an enumerate, a description list with bold run-in term
labels, low-level sectioning down to subsubsection depth, paragraph and
subparagraph headings followed by real text, inline mathematics embedded inside
list items, and cross-references that resolve both to sections and to an
individually labelled enumerate item. The prose is substantive so that a
text-difference metric over the body remains informative.
\end{abstract}

\section{Review Protocol}
\label{sec:protocol}

The review protocol is organised as a hierarchy of stages, sub-tasks, and
checks. The outermost layer is an ordered list of stages; within a stage the
operative sub-tasks form an unordered list; and within a sub-task the concrete
checks are again ordered. This arrangement produces three levels of nesting and
is the central structure under test in \cref{sec:protocol}.

\begin{enumerate}[label=\arabic*.]
\item \textbf{Intake.} The reviewer confirms that each record is admissible
  before any outcome data are inspected. The admissible sample size is
  $n_0 \leq n$, where $n$ is the raw count of submitted records.
  \begin{itemize}
  \item Structural fields are validated against the intake schema:
    \begin{enumerate}[label=(\roman*)]
    \item the site identifier is present and resolves to a known site;
    \item the submission timestamp lies inside the collection window
      $[t_0, t_1]$; \label{item:window}
    \item the record carries a reviewer initial recognised by the roster.
    \end{enumerate}
  \item Content fields are screened for completeness, with any field whose
    missingness exceeds the cutoff $m^{\ast} = 0.10$ flagged for a second pass.
  \end{itemize}
\item \textbf{Adjudication.} Flagged records are escalated and resolved under a
  fixed decision rule.
  \begin{itemize}
  \item Duplicate initials are reconciled against the supervisor code before the
    record advances.
  \item Borderline eligibility is decided by majority among three readers, so a
    record is retained when at least $\lceil 3/2 \rceil = 2$ readers agree.
  \end{itemize}
\item \textbf{Sign-off.} The final analytic status is recorded, after which the
  record is locked and the outcome fields become readable.
\end{enumerate}

The timestamp condition in \cref{item:window} is the single most common cause of
rejection at intake, because submissions assembled before the collection window
opens are silently discarded by the upstream system.

\section{Definitions and Roles}
\label{sec:definitions}

The following description list assigns a bold run-in label to each defined term.
Run-in labels are a frequent point of failure in conversion because the label
must remain on the same line as the explanatory text that follows it.

\begin{description}
\item[Reviewer] An individual who performs the intake and adjudication tasks of
  \cref{sec:protocol} and whose initials appear in the audit log.
\item[Supervisor] A reviewer with the additional authority to reconcile
  duplicate initials and to break ties during adjudication.
\item[Collection window] The closed interval $[t_0, t_1]$ during which a
  submission is considered timely; submissions outside the interval are
  inadmissible at intake.
\item[Analytic status] The terminal label assigned at sign-off, which fixes
  whether a record enters the analysis sample.
\end{description}

\subsection{Operational Notes}
\label{sec:opnotes}

This subsection records operational details that elaborate the roles defined in
\cref{sec:definitions}. The notes are written as ordinary prose so that
sectioning depth is exercised without relying on further lists.

\subsubsection{Escalation Path}
\label{sec:escalation}

When a record cannot be resolved by the assigned reviewer, it follows a fixed
escalation path to the supervisor. The escalation is logged with a timestamp,
and the supervisor's resolution overrides the original disposition. This
subsubsection establishes the deepest standard sectioning level in the document,
and \cref{sec:escalation} is referenced again in the closing section to confirm
that low-level section labels survive conversion.

\paragraph{Logging convention.}
Every escalation writes one line to the audit log. The line contains the record
identifier, the originating reviewer, the supervisor, and the resolution code,
in that fixed order. A paragraph heading is used here, with substantive text
immediately following the heading on the same logical block, because run-in
paragraph titles are a known source of conversion drift.

\subparagraph{Retention rule.}
A logged escalation is retained for the full retention horizon $H$ even after the
parent record is locked, so that the resolution remains auditable. The
subparagraph heading is the lowest heading level in the document and is included
specifically to test whether the converter preserves the distinction between a
paragraph and a subparagraph.

\section{Worked Counts}
\label{sec:counts}

The protocol of \cref{sec:protocol} produces several counts that can be checked
arithmetically. The list below embeds inline mathematics inside each item so
that short formulas must survive within list structures.

\begin{enumerate}[label=(\alph*)]
\item The retained fraction at intake is
  $r = n_0 / n$, which lies in the unit interval $0 \leq r \leq 1$.
\item A record advances past adjudication when at least two of three readers
  agree, an event whose count is $k \in \{2, 3\}$.
\item The number of distinct sites observed is
  $s = \lvert \{\, \text{site}(i) : i \in \text{sample} \,\} \rvert$, which is
  bounded above by the number of registered sites.
\end{enumerate}

These quantities connect the structural rules of \cref{sec:protocol} to the
definitions of \cref{sec:definitions}, and the escalation logging of
\cref{sec:escalation} ensures that every adjudication decision underlying item
(b) is recoverable from the audit log. The timestamp gate first introduced in
\cref{item:window} therefore propagates through the entire pipeline, from intake
counts through to the retained analytic sample.

\end{document}
